\subsection{Limitaciones del Estudio y Restricciones Técnicas}
\label{subsec:limitaciones_estudio}

A pesar del rigor metodológico aplicado, el presente trabajo presenta limitaciones técnicas que deben ser consideradas:

\begin{itemize}
    \item \textbf{Dependencia de API Propietaria en la Etapa 2 de Reordenación:} Aunque la Etapa 1 de recuperación densa soporta ya ejecución íntegramente local mediante \textit{Qwen3-VL-Embedding-2B} y \textit{EmbeddingGemma 2} (Subsección~\ref{subsec:comparativa_backends_locales}), el verificador bimodal de segunda etapa continúa dependiendo del modelo remoto \textit{Gemini 3.5 Flash Lite}. Aunque se implementaron reintentos exponenciales y cortafuegos de cuota ($1.000\text{ RPD}$), un sistema 100\% desconectado de Internet requiere incorporar también un reordenador local.
    \item \textbf{Precisión Fonética del Modelo Whisper \texttt{base}:} Para hacer viable la ejecución local en tiempo real sobre CPU, se seleccionó el modelo Whisper en su variante \texttt{base}. Aunque demostró un rendimiento notable, ocasionalmente comete errores de transcripción en palabras técnicas complejas o modismos coloquiales muy específicos, los cuales se propagan al texto inyectado en el vectorizador.
    \item \textbf{Sensibilidad del Algoritmo SSIM ante Cambios Graduales:} La detección de cortes basada exclusivamente en SSIM es idónea para transiciones de plano netas, pero muestra menor sensibilidad ante transiciones lentas por disolución encadenada (\textit{cross-dissolve}) o movimientos panorámicos continuos de cámara, dependiendo en esos casos de la cota máxima temporal ($T_{\max} = 8,0\text{ s}$ por defecto y $5,0\text{ s}$ en el banco de código estático).
    \item \textbf{Tamaño del Corpus de Validación:} Si bien el banco de pruebas (164 fragmentos, 138 consultas) fue anotado minuciosamente a mano para garantizar un \textit{ground truth} de máxima calidad, la escala total del metraje (aproximadamente 8,1 minutos) representa una prueba de concepto que debe validarse a futuro sobre conjuntos de datos públicos masivos de cientos de horas de duración (como \textit{Video-MME} o \textit{LongVideoBench}) \cite{video_rag2024, rethinking_long_video2026}.
\end{itemize}
